\documentclass{article}
\usepackage[utf8]{inputenc}
\usepackage[american]{babel}
\usepackage{tikz}
\usepackage{standalone}
\usepackage{parskip}
\usepackage{float}
\title{Architecting for Action: \\ \large Data Lineage and Analysis via the T-DAR Framework}
\author{Tim Didelot \\ \texttt{contact@timdidelot.fr}}
\date{\today}

\begin{document}
\maketitle

\begin{abstract}
\noindent
Data lineage is a critical aspect of modern data management, providing insights into the flow and transformation of data across complex ecosystems. This paper introduces the T-DAR framework, a novel approach to analyzing data lineage that emphasizes the relationship between targeted observations, data assets, and actionable results. By leveraging the T-DAR framework, organizations can enhance their understanding of data processes, improve decision-making, and ensure compliance with regulatory requirements.
\end{abstract}

\section{Introduction}
In the era of Big Data, organizations have adopted a data-driven approach to decision-making. However, the complexity of data ecosystems often leads to challenges in understanding the flow and transformation of data. This paper introduces the T-DAR framework, which stands for Targeted Data-Actions-Results, as a methodology for analyzing data lineage and ensuring actionable insights.

\section{T-DAR Framework Overview}
The T-DAR framework is inspired by the Data-Action-Result (DAR) model of most healthcare information systems.
\begin{figure}[H]
    \centering
    \resizebox{0.85\textwidth}{!}{
        \makebox[\textwidth][c]{
            \input{out/uml/tdar_architecture/tdar_architecture.latex}
        }
    }
    \caption{T-DAR Architecture Overview}
    \label{fig:tdar_architecture}
\end{figure}

As illustrated in Figure \ref{fig:tdar_architecture}, the architecture of the T-DAR framework is centered around a \textit{Targeted Observation}, which is directly linked to a specific \textit{Data Asset} within the ecosystem. Rather than treating data lineage as a static representation of data streams, the T-DAR framework takes advantage of the dynamic nature of data environments, allowing for real-time analysis and decision-making. It provides a way to analyse not only results, but also the medium through which the results are obtained, and the actions that lead to those results.

To operationalize this dynamic analysis, every \textit{Trgeted Observation} acts as a container for a chronological sequence of \textit{DAREntry} records. These entries are categorized into three distinct phases: \textit{Data}, \textit{Action}, and \textit{Result}, which together from the foundation of the framework:
\begin{itemize}
    \item \textbf{Data (D) -- The Assessment:} A Data entry represents a factual snapshot, an audit trail, or a metric collected from the \textit{Data Asset}.
    \item \textbf{Action (A) -- The Intervention:} An Action entry represents a specific action taken in response to the data or the results of the data.
    \item \textbf{Result (R) -- The Outcome:} A Result entry represents the outcome or consequence of the preceding action taken.
\end{itemize}

\end{document}